\begin{section}{The Extended Real Number System}

The Extended real number system $\overline{\R}$ is defined as
\[
\overline{\R} \coloneqq \R \cup \set{-\infty, +\infty}.
\]
We define the following operations on $\overline{\R}$:
\begin{itemize}
    \item $x + (+\infty) \coloneqq +\infty$ for all $x \in \R \cup \set{+\infty}$.
    \item $x + (-\infty) \coloneqq -\infty$ for all $x \in \R \cup \set{-\infty}$.
    \item $x \cdot (\pm \infty) \coloneqq \pm \infty$ for all $x \in \R^{+} \cup \set{+\infty}$.
    \item $x \cdot (\pm \infty) \coloneqq \mp \infty$ for all $x \in \R^{-} \cup \set{-\infty}$.
    \item $\dfrac{x}{\pm \infty} \coloneqq 0$ for all $x \in \R$.
\end{itemize}

$\overline{\R}$ is not a field, but it is an ordered set with the following order relation:
\[
-\infty < x < +\infty \quad \text{for all } x \in \R.
\]

\bigskip

\begin{theorem} \label{thm:SupremumInfimumInExtendedRealNumberSystem}
    Any subset of $\overline{\R}$ has a least upper bound and a greatest lower bound in $\overline{\R}$.
\end{theorem}

\begin{proof}
    Let $E \subseteq \overline{\R}$. If $E = \emptyset$, then $\sup E = -\infty$ and $\inf E = +\infty$.
    If $E \neq \emptyset$, then we have the following cases:
    \begin{itemize}
        \item If $E$ is bounded above in $\R$, then $\sup E$ exists in $\R$ by the least upper bound property of $\R$.
        \item If $E$ is not bounded above in $\R$, then $\sup E = +\infty$.
        \item If $E$ is bounded below in $\R$, then $\inf E$ exists in $\R$ by the greatest lower bound property of $\R$.
        \item If $E$ is not bounded below in $\R$, then $\inf E = -\infty$.
    \end{itemize}
\end{proof}

\end{section}